\section{模型假设与符号}

\subsection{模型假设}

\begin{enumerate}[leftmargin=1.9em,itemsep=1pt]
\item[A1] \textbf{目标已知在圆柱内.} 搜救开始时，搜救方知道目标在水平半径 100 m、
      垂直 $\pm$10 m 的圆柱内；圆柱中心由上层（题设分布）给定。
\item[A2] \textbf{初始位置分布.} 失事海域为 4$\times$4 km 矩形，水平均匀；
      失事深度服从 $p(z)\propto e^{-(z-200)/200}$，$z\ge 200$ m，且不超过当地海底
      深度与额定潜深 1000 m。
\item[A3] \textbf{失效树完备.} 5 种失效模式（S0--S4）覆盖主要情形，模式概率由专家判断给出，
      模式间互斥且完备（见 \tabref{tab:failure}）。
\item[A4] \textbf{艇体可等效为球.} 用同体积球体近似耐压壳的阻力与附加质量，
      附加质量系数取球体理论值 0.5。
\item[A5] \textbf{浮力异常为常数.} 失事后艇体以新的平均密度
      $\rho_b=\rho(z_0)+\Delta\rho$ 运动，$\Delta\rho$ 由压载状态决定，不随时间变化。
\item[A6] \textbf{流场为预报场.} 用 HYCOM 的四维流场（0--96 h）作为真值；搜索计划
      与真值使用同一流场（“完美预报”），该假设的影响在 §\ref{sec:carib} 讨论。
\item[A7] \textbf{海面为吸收边界.} 到达水面的目标持续漂浮并随表层流与风漂移；
      到达海底的目标静止，或在坡度超过静摩擦角时沿坡滑移。
\item[A8] \textbf{压溃不参与模拟.} 超过额定潜深的目标在模拟中继续下沉；
      “压溃”作为\emph{事后统计}给出（给定假设的压溃深度，给出超深概率曲线）。
\item[A9] \textbf{探测相互独立.} 各装备、各次搜索对同一目标的探测事件相互独立，
      多轮搜索的探测概率按 $1-\prod(1-p_i)$ 累积。
\item[A10] \textbf{目标类别三分.} 按目标当前位置分为水面（$z\le2$ m）、水柱
      （$2<z<$底−6 m）、海底（$z\ge$底−6 m）三类，装备对不同类别有不同扫宽。
\item[A11] \textbf{装备可用性.} 单船方案中，除雷达外的装备从出发岛出发；
      多船方案中 t=0 只有母船，增援船与 AUV 在 3 h 后到位（可做敏感性）。
\item[A12] \textbf{成本为量级估计.} 装备采购/动员/使用成本取公开市场量级估计，
      在报告中标注为“估计值”并做敏感性讨论。
\end{enumerate}

\subsection{符号表}

\begin{table}[H]
\centering\small
\caption{主要符号及其含义（单位）}
\label{tab:symbols}
\begin{tabular}{@{}llll@{}}
\toprule
符号 & 含义（单位） & 符号 & 含义（单位）\\
\midrule
$m$ & 潜水器质量（kg） & $m^*$ & 附加质量后的有效质量（kg）\\
$V$ & 排水体积（m$^3$） & $A$ & 等效迎流面积（m$^2$）\\
$\rho(z)$ & 海水现场密度（kg/m$^3$） & $\rho_b$ & 艇体平均密度（kg/m$^3$）\\
$\Delta\rho$ & 浮力异常 $\rho_b-\rho(z_0)$（kg/m$^3$） & $z_0$ & 失事深度（m）\\
$w$ & 垂向速度，向下为正（m/s） & $w_t$ & 终端速度（m/s）\\
$C_D$ & 阻力系数（无量纲） & $g$ & 重力加速度（m/s$^2$）\\
$\mu_s,\mu_k$ & 静/动摩擦系数（无量纲） & $\theta$ & 海底坡度（°）\\
$\bm u=(u,v)$ & 水平流速（m/s） & $\bm u_b$ & 系统偏差（m/s）\\
$\varepsilon$ & Ornstein--Uhlenbeck 流速扰动（m/s） & $\tau_{\rm OU}$ & 扰动相关时间（h）\\
$K_h$ & 水平湍流扩散系数（m$^2$/s） & $\sigma_0$ & 初始定位误差尺度（m）\\
$W_c$ & 对类别 $c$ 的有效扫宽（m） & $v$ & 搜索速度（m/s）\\
$\tau$ & 单格停留时间（h） & $A_c$ & 网格面积（km$^2$）\\
$S_i$ & 粒子 $i$ 尚未被发现的概率 & $\beta$ & 重复扫描惩罚系数\\
$A_{50},A_{95}$ & 50\%/95\% 包含域面积（km$^2$） & $P_{72}$ & 72 h 内发现概率\\
\bottomrule
\end{tabular}
\end{table}

表中 $A_{50}$ 与 $A_{95}$ 的含义是：把位置概率从高到低累加，覆盖总概率 50\%/95\%
所需的最小面积。它直接对应“搜救方需要搜索多大范围”，是定位模型与搜索模型之间的桥梁。
